\medskip\noindent\textbf{Supplementary Note 5. Matched batch removal $b^*$ and the low-signal rule.}\par\smallskip
\begingroup\small\setlength{\parskip}{3pt}
Penalty strengths multiply losses with different scales, so equal strengths do not imply equal batch removal. We therefore compare biology at a common batch score, $b^*$. Let $b_0$ be plain scVI's mean batch score over seeds 100--102. For each penalty, take its highest seed-mean batch score among strengths with no failed fit. The smallest of these maxima is $b_{\mathrm{common}}$, the highest batch score all penalties reach. Then
\begin{equation}
b^*=b_0+\tfrac12\,(b_{\mathrm{common}}-b_0).
\label{eq:s5-bstar}
\end{equation}
Each curve starts at plain scVI and joins strengths in increasing order, using mean scores over successful fits. Strengths with no successful fits are omitted. Biology is read at the first point equal to $b^*$ within $10^{-9}$, or linearly interpolated in batch score across the first bracket. Later crossings are ignored. If plain scVI's batch score is at or above $b^*$, the curve gives no reading. Every full-seed curve crosses when $b_{\mathrm{common}}>b_0$, but curves with a seed omitted or with new seeds need not cross the fixed target.

A pair is low-signal when $b_{\mathrm{common}}-b_0<0.02$, placing $b^*$ less than 0.01 above plain scVI. Three pairs meet this rule, all with the batch-conditioned decoder: immune has scores 0.638 and 0.657, gap 0.019; simulation 1 has 0.812 and 0.819, gap 0.007; simulation 2 has 0.656 and 0.665, gap 0.009. No unconditioned pair is low-signal.

Low-signal pairs are excluded from rankings, their uncertainty summaries, hypervolume and main-text label-transfer comparisons. Their readings remain in Supplementary Tables S6, S8 and S11. Non-dominated shares, collapse rates, single-strength placement and the direct biology check retain them. All 96 full-seed curves reached $b^*$, including the low-signal pairs.

Label transfer and the direct biology check require one fitted embedding. Among failure-free strengths bracketing the first crossing, they use the strength whose mean batch score is nearer $b^*$, breaking ties towards the smaller strength. If the lower point is plain scVI, they use the upper strength. Edge and extension strengths are eligible. Batch scores alone choose this matched strength, and its distance from $b^*$ is reported. Supplementary Note 10 describes the change from the earlier selection rule.
\par\endgroup
